\subsection{操作のつながりを確かめる接続演習}\label{節：接続演習}
  新しい内容を,既習の操作へ戻して確認する7題である.
  問1・2では置き換える量を選び,問3・4では更新の範囲を組み替え,
  問5〜7では周期・近似・和を調べる.
  解法名だけでなく,選んだ操作で何が簡単になったかを述べよう.

  \subsubsection{問題}
  \begin{enumerate}[label=\textbf{問\arabic*.},format=\bookitemlink{practice-connections}{problem}{answer}]
    \item $a_1=3$,
      $a_{n+1}=\frac{n+2}{n}a_n+n(n+1)(n+2)$ $(n\ge1)$とする.
      $b_n=\frac{a_n}{n(n+1)}$の更新を求め,\bookterm{general}{一般項}を求めよ.
      また,同じ更新をする一つの\bookterm{sequence}{数列}$v_n$を見つけ,
      $a_n-v_n$を追う別解を説明せよ.
    \item $a_1=1$, $a_2=2$,
      $a_{n+2}=\frac{a_{n+1}^2+1}{a_n}$ $(n\ge1)$とする.
      各項が正であることを示し,
      $I_n=\frac{a_n^2+a_{n+1}^2+1}{a_na_{n+1}}$を使って\bookterm{general}{一般項}を求めよ.
      \bookterm{invariant}{不変量}の値は\bookterm{initial}{初期条件}を変えるとどう変わるか.
    \item $n,m\ge0$とする.$A_{n,0}=A_{0,m}=1$,
      $A_{n,m}=A_{n-1,m}+A_{n,m-1}$ $(n,m\ge1)$から,
      $A_{n,m}=\binom{n+m}{m}$を$n+m$に関する帰納法で証明せよ.
      また右と上に進む格子経路の数として説明せよ.
    \item $a_1=2$とし,
      \[
        a_{n+1}=\begin{cases}
          -a_n&(n\text{が奇数}),\\
          2a_n+1&(n\text{が偶数})
        \end{cases}
      \]
      とする.奇数番目と偶数番目の\bookterm{general}{一般項}を求めよ.
      係数が周期的であることと,\bookterm{sequence}{数列}自身が周期的であることは同じか.
    \item $a_1=1$, $a_2=2$,
      $a_{n+2}=\frac{1+a_{n+1}}{a_n}$ $(n\ge1)$とする.
      周期を使って\bookterm{general}{一般項}を定め,最小周期を求めよ.
      正の初期値を変えて定\bookterm{sequence}{数列}となる場合も求めよ.
    \item $x_1=3$, $y_1=2$,
      $x_{k+1}=3x_k+4y_k$, $y_{k+1}=2x_k+3y_k$ $(k\ge1)$とする.
      $x_k^2-2y_k^2=1$を証明せよ.
      $r_k=x_k/y_k$について,$r_k-\sqrt2$を$y_k,r_k$で表して\bookterm{limit}{極限}を求め,
      $\frac12(r_k+2/r_k)=r_{2k}$も示せ.
    \item $j^3=6\binom j3+6\binom j2+\binom j1$を確かめよ.
      二項係数の\bookterm{difference-operator}{差分}を使って,$\sum_{j=0}^{N-1}j^3$ $(N\ge0)$を求めよ.
      単項式の係数比較で求める方法と,何が違うか述べよ.
  \end{enumerate}

  \subsubsection{方針}
  \begin{enumerate}[format=\bookitemlink{practice-connections}{hint}{problem}]
    \item 割る因子の隣り合う比を確認する.残った\bookterm{difference}{階差}は$n$となる.
      \bookterm{normalize}{正規化}後の一つの解から,引く量$v_n$を逆に作る.
    \item 正値性を先に帰納的に確かめる.$xz=y^2+1$から$I_{n+1}=I_n$を示し,
      初期値で係数を決める.
    \item 内側の二つの位置では添字の和が一つ小さい.
      境界を先に確かめ,二項係数の加法公式を使う.
    \item 二回の更新をまとめる.$x_k=a_{2k-1}$は\bookterm{constant}{定数係数}の一次\bookterm{recurrence}{漸化式}になる.
    \item 最初の七項を計算し,直前の二項の組が戻ることを確かめる.
      \bookterm{initial}{初期条件}が変わる場合は文字で計算する.
    \item 保存される式は更新の前後で展開して比べる.
      誤差は$x_k^2-2y_k^2$の因数分解,添字の倍増は二次式の平方から得る.
    \item 連続整数の積へ展開して恒等式を確かめる.
      各二項係数の和は,上の添字$N$で下の添字を一つ増やしたものになる.
  \end{enumerate}

  \subsubsection{解答}
  \begin{enumerate}[format=\bookitemlink{practice-connections}{answer}{problem}]
    \item $b_{n+1}-b_n=n$, $b_1=\frac32$より,
      \[
        a_n=\frac{n(n+1)(n^2-n+3)}2.
      \]
      $v_n=\frac{n^2(n-1)(n+1)}2$を選ぶと,
      $a_n-v_n=\frac32n(n+1)$となる.
    \item $I_n=3$なので$a_{n+2}=3a_{n+1}-a_n$となる.
      $\alpha=\frac{3+\sqrt5}{2}$, $\beta=\frac{3-\sqrt5}{2}$とすると,
      \[
        a_n=\frac{(2-\beta)\alpha^{n-1}+(\alpha-2)\beta^{n-1}}{\sqrt5}.
      \]
      正の初期値$x,y$では$I_n=(x^2+y^2+1)/(xy)$となる.
    \item 境界と加法公式により$A_{n,m}=\binom{n+m}{m}$である.
      $n+m$歩から上への$m$歩を選ぶ数でもある.
    \item $k\ge1$について,
      \BookMathLayout{\[
        a_{2k-1}=\frac{1+5(-2)^{k-1}}3,\qquad
        a_{2k}=-\frac{1+5(-2)^{k-1}}3.
      \]}{\[
    \begin{aligned}
      a_{2k-1}&=\frac{1+5(-2)^{k-1}}3,\\
      a_{2k}&=-\frac{1+5(-2)^{k-1}}3.
    \end{aligned}
  \]}
      絶対値が非有界になるので,\bookterm{sequence}{数列}自身は周期的ではない.
    \item 各項は$1,2,3,2,1$を繰り返す.
      $n-1$を$5$で割った余りが$0,1,2,3,4$のとき,
      $a_n$はそれぞれ$1,2,3,2,1$である.最小周期は$5$.
      正の定\bookterm{sequence}{数列}は$a_1=a_2=(1+\sqrt5)/2$の場合である.
    \item $x_k^2-2y_k^2=1$が保たれ,
      \[
        r_k-\sqrt2=\frac1{y_k^2(r_k+\sqrt2)},\qquad
        \lim_{k\to\infty}r_k=\sqrt2.
      \]
      また$r_{2k}=\frac12(r_k+2/r_k)$である.
    \item
      \BookMathLayout{\[
        \sum_{j=0}^{N-1}j^3
          =6\binom N4+6\binom N3+\binom N2
          =\frac{N^2(N-1)^2}{4}.
      \]}{\[
    \begin{aligned}
      \sum_{j=0}^{N-1}j^3&=6\binom N4+6\binom N3+\binom N2\\
      &=\frac{N^2(N-1)^2}{4}.
    \end{aligned}
  \]}
      \bookterm{difference-operator}{差分}に適した部品へ表し直してから,部品ごとに和を取った.
  \end{enumerate}

  \subsubsection{解法}
  \begin{enumerate}[format=\bookitemlink{practice-connections}{solution}{problem}]
    \item \textbf{倍率をそろえる.}
      $h_n=n(n+1)$は$h_{n+1}/h_n=(n+2)/n$を満たす.
      両辺を$(n+1)(n+2)$で割ると$b_{n+1}=b_n+n$であり,
      $b_n=\frac32+\sum_{j=1}^{n-1}j=\frac{n^2-n+3}{2}$となる.
      \par\smallskip\textbf{引く別解との比較.}
      \bookterm{normalize}{正規化}後の式では,$w_n=n(n-1)/2$も$w_{n+1}-w_n=n$を満たす.
      $v_n=h_nw_n$は元の\bookterm{inhomogeneous}{非同次}の更新を満たすため,
      $d_n=a_n-v_n$について$d_{n+1}=\frac{n+2}{n}d_n$となる.
      $d_1=3$から$d_n=\frac32n(n+1)$を得る.
      この場合は割る方法から引く量も発見できた.
      \sectionref*{節：変数係数の一次漸化式}へ戻って比べよう.
    \item \textbf{正値性と保存を確認する.}
      正の二項から次も正になる.$x=a_n,y=a_{n+1},z=a_{n+2}$とすると,
      $xz=y^2+1$から$I_n=I_{n+1}=(x+z)/y$である.
      $I_1=(1+4+1)/2=3$なので$x+z=3y$となる.
      二つの\bookterm{root}{特性根}で$a_n=C\alpha^{n-1}+D\beta^{n-1}$と置き,
      $C+D=1$, $C\alpha+D\beta=2$を解くと解答の係数になる.
      $I_n$は初期値全体で共通の定数ではなく,選んだ\bookterm{sequence}{数列}の中で一定である.
      \sectionref*{節：不変量から線形の関係}へ.
    \item \textbf{境界から内側へ証明する.}
      $n=0$または$m=0$では式は$1$になる.
      $n,m\ge1$の場合,$n+m$より小さい添字の和で式が成立すると仮定し,
      \BookMathLayout{\[
        A_{n,m}=\binom{n+m-1}{m}+\binom{n+m-1}{m-1}
               =\binom{n+m}{m}
      \]}{\[
    \begin{aligned}
      A_{n,m}&=\binom{n+m-1}{m}+\binom{n+m-1}{m-1}\\
      &=\binom{n+m}{m}
    \end{aligned}
  \]}
      を得る.格子経路では最後の一歩が右か上かで分けた和であり,
      全歩の並べ方を直接数えても同じ二項係数になる.
      \sectionref*{節：二重数列と二方向の階差}へ.
    \item \textbf{一周期分をまとめる.}
      $x_k=a_{2k-1}$とすると,$a_{2k}=-x_k$から
      $x_{k+1}=-2x_k+1$となる.
      \bookterm{fixed}{不動点}$1/3$を引けば$x_{k+1}-1/3=-2(x_k-1/3)$であり,
      $x_1=2$から解答の式が得られる.
      係数の周期が$2$でも,一周期後の倍率$-2$によって項の絶対値は大きくなる.
      \sectionref*{節：周期的な係数}へ.
    \item \textbf{二項の組が戻る.}
      $a_3=3,a_4=2,a_5=1,a_6=1,a_7=2$なので,
      $(a_6,a_7)=(a_1,a_2)$である.更新の一意性から全体が繰り返す.
      周期$5$の最小周期は$1$または$5$であり,定\bookterm{sequence}{数列}ではないので$5$である.
      定数$t>0$では$t^2=t+1$が必要十分である.
      \sectionref*{節：非線形の周期性}へ.
    \item \textbf{保存・近似・添字の倍増を結ぶ.}
      更新を代入して展開すれば
      \[
        (3x+4y)^2-2(2x+3y)^2=x^2-2y^2
      \]
      となり,初期値での値$1$が保存される.
      正値性と$y_{k+1}>3y_k$から$y_k\to\infty$であり,
      $(r_k-\sqrt2)(r_k+\sqrt2)=1/y_k^2$から誤差と\bookterm{limit}{極限}が得られる.
      また更新と初期値から$x_k+y_k\sqrt2=(3+2\sqrt2)^k$である.
      平方して$x_{2k}=x_k^2+2y_k^2$, $y_{2k}=2x_ky_k$を得れば,
      比は\bookterm{newton}{ニュートン法}の式になる.
      \sectionref*{節：ペル方程式と近似分数}へ.
    \item \textbf{操作が簡単になる表示を選ぶ.}
      $j(j-1)(j-2)+3j(j-1)+j=j^3$だから,提示された恒等式が成り立つ.
      \bookterm{difference-operator}{差分}の和の公式をそれぞれに使うと,
      $6\binom N4+6\binom N3+\binom N2$となる.
      $N\ge4$では階乗表示を整理して解答の式を得る.
      $N=0,1,2,3$でも,範囲外の二項係数を零とする規約で直接確かめられる.
      係数比較で差を逆に戻す方法に対し,ここでは\bookterm{difference-operator}{差分}で添字が一つ下がる部品を使った.
      \sectionref*{節：二項係数と差分}へ.
  \end{enumerate}
